\section{Limits and Openings}
\label{sec:limites_ouvertures}

\subsection{What is solidly established}
\label{sec:acquis}

The HP-PT programme, in its current state, provides:
\begin{itemize}[nosep,leftmargin=*]
\item the closed-form expression of the amplitudes $\kappap(s)$ and
the exact Fredholm representation of $\Rres(s)$
(thm.~\ref{thm:fredholm});
\item the canalized operator $\Hilb_{\mathrm{PT}}$ for the coupling
with canonical $\epsstar^{\mathrm{geom}}$ without free parameter
(lem.~\ref{lem:epsstar_geom});
\item the dynamical cut-off $\umax(\gamma) = \gamma/\sqrt{2\pi}$
forced by the symplectic Planck cell
(prop.~\ref{prop:cellule_planck});
\item the exact semi-classical density
$\NPT(\gamma) = N_{\mathrm{RvM}}(\gamma) + 1/8$ (identical functional
form) (prop.~\ref{prop:densite_NPT});
\item the defect indices $(1, 1)$ and the antiperiodic boundary
condition forced by $T_3$ (prop.~\ref{prop:antiperiodicite});
\item the exact analytic trace formula $-d \log \Rres/ds$
(thm.~\ref{thm:trace_formula});
\item the operatorial regularised trace $\Treg(s)$ with
distributional poles at the zeros of $\Rres$
(\S\ref{sec:trace_reg});
\item the triple cohomological identification of the residue $1/8$
(\eqref{eq:triple_identification});
\item the structural Higgs $\leftrightarrow \zeta$ duality
(conjecture~\ref{conj:K4}).
\end{itemize}

\paragraph{Summary table of epistemic statuses.}

\begin{longtable}{@{}p{8cm}p{2.5cm}p{4cm}@{}}
\toprule
\textbf{Result} & \textbf{Status} & \textbf{Source}\\
\midrule
\endhead
$\sper = 1/2$ ($\Tone$) & \textbf{[Thm]} & \cite{PT_M1}\\
$\det(I - \DRop)^{-1} = \Rres(s)$, $\Reop(s) > 1$ & \textbf{[Thm]} & thm.~\ref{thm:fredholm}\\
Trace formula $-d\log \Rres/ds$ in $\Reop(s) > 1$ & \textbf{[Thm]} & thm.~\ref{thm:trace_formula}\\
Defect indices $(1,1)$ and AP BC $\theta = \pi$ & \textbf{[Thm]} & prop.~\ref{prop:antiperiodicite}\\
% ERRATUM E-PM-8 (2026-08-15): source made precise. The [Thm] covers ONLY the
% algebraic half $c_1 = 1/2$ of thm:berry_3pi; the holonomy half ($\exp(3\pi i)$)
% of that same theorem is [COND] (see \S 5). Status unchanged.
$c_1 = 1/2$ on spinor bundle $\qplus/\qminus$ & \textbf{[Thm]} & thm.~\ref{thm:berry_3pi} (algebraic half)\\
Density $\NPT = N_{\mathrm{RvM}} + 1/8$ (asymptotic identity) & \textbf{[Der]} & prop.~\ref{prop:densite_NPT}\\
$1/8 = c_1/N_{\mathrm{corners}}$ (strict APS rigour) & \textbf{[Thm partial]} & prop.~\ref{prop:c1_sur_N}\\
$\zetaplus\,\zetaminus \ne 0$ on critical line & \textbf{[Conj]} & rem.~\ref{rem:non_annulation}\\
Strict analytic continuation of $\Rres(s)$ outside $\Reop(s) > 1$ & \textbf{[Open]} & op.~1 below\\
Distributional poles of $\Treg(s)$ on $\Reop(s) = 1/2$ only & \textbf{[Open]} ($\equiv$~RH) & op.~2 below\\
Canonical choice of $u_\star$ ($\sqrt{2\pi}$ vs $\log 3$) & \textbf{[Open]} & op.~4 below\\
Structural Higgs $\leftrightarrow \zeta$ duality (K4) & \textbf{[Strong Conj]} & conj.~\ref{conj:K4}\\
$\Delta_N(L \bmod q) = 1/(8q)$ for Dirichlet (P3) & \textbf{[Falsified, $T \le 200$]} & \S\ref{sec:M1bis}\\
$G_{\mathrm{HP4.d}}$: direct spectrum = $\{\gamma_n\}$ & \textbf{[Falsified]} & \S\ref{sec:trace_reg}\\
\bottomrule
\caption{Epistemic statuses of the programme's results.}
\label{tab:statuts}
\end{longtable}

\subsection{What remains open}
\label{sec:ouvertures}

Four structural openings, plus an underlying non-vanishing
hypothesis.

\paragraph{Opening 1 (G\_HP3.a).}
Analytic continuation of $\Rres(s)$ outside $\Reop(s) > 1$. The trace
$\Treg(s)$ is defined by series in $\Reop(s) > 1$ and then extended
distributionally to $\Reop(s) = 1/2$; the intermediate region
$1/2 < \Reop(s) < 1$ requires classical analysis of $\Rres$. This is
essentially the classical problem of $\zeta$ continuation.

\paragraph{Opening 2 (strict RH).}
Prove that \textbf{all} distributional poles of $\Treg(s)$ lie on
$\Reop(s) = 1/2$. The programme establishes that the structure
favours this line (four equivalent mechanisms,
\S\ref{sec:localisation_critique}), but does not prove the strict
exclusion of off-line poles. This is essentially RH.

\paragraph{Opening 3 (strict APS K3).}
Formal Atiyah--Patodi--Singer rigour for the cornered manifold of
the BK double barrier, beyond the area calculation and numerical
verification. The cohomological mechanism is identified
(\S\ref{sec:indice_APS}), its strict formalisation remains to be
developed.

\paragraph{Opening 4 (K1, canonical choice of $u_\star$).}
The lower barrier $u_\star$ admits two natural candidates:
$u_\star = \sqrt{2\pi}$ (fixed point of the involution
$u \leftrightarrow p_u$, chosen here) and $u_\star = \log 3$
(first PT active prime, arithmetic choice). Both produce the same
asymptotic density $N(\gamma)$ but different $O(1/\gamma)$
corrections. The canonical status of the choice between them remains
open. The article retains $u_\star = \sqrt{2\pi}$ by involution
symmetry, without claiming to close K1.

\paragraph{Underlying \textbf{[Conj]} hypothesis: non-vanishing $\zetaplus\,\zetaminus \ne 0$.}
The identification between zeros of $\Rres$ and zeros of $\zeta$
used throughout the article relies on this non-vanishing. It is
verified numerically on the tested zones but not proved
(rem.~\ref{rem:non_annulation}). A partial failure would introduce
extra zeros (zeros of $\zeta_\pm$) which would translate into extra
poles of $\Treg(s)$; these possible zeros must be treated as
conjecturally absent.

\subsection{Canonical position of the programme}
\label{sec:position_canonique}

The HP-PT programme is a precise operatorial framework for the
zeros of $\zeta$. It provides:
\begin{enumerate}[nosep,leftmargin=*]
\item a closed-form expression for the zeros via $\kappap(s)$
(intermediary $\DRop$);
\item a canonical regularised trace formula $\Treg(s)$;
\item four equivalent mechanisms for $\Reop(s) = 1/2$;
\item a continuous/discrete dictionary unified by persistence;
\item a structural explanation of the known numerical coincidences
($2\pi$, $7/8$, $1/8$, critical line).
\end{enumerate}

It does not bring:
\begin{enumerate}[nosep,leftmargin=*]
\item a strict proof of classical RH;
\item a shortcut to the distributional analysis of the regularised
trace.
\end{enumerate}

This distinction is maintained throughout the programme: one seeks
to \textbf{understand} the zero structure, not to \textbf{prove}
them on the critical line. The passage to formal proof requires
fine Schwartz analysis on $\Treg(s)$, which is no easier than the
classical RH proof.

\subsection{Comparison with Berry--Keating--Connes}
\label{sec:comparaison_BKC}

\begin{table}[ht]
\centering
\small
\begin{tabular}{@{}p{3cm}p{4.5cm}p{5cm}@{}}
\toprule
\textbf{Aspect} & \textbf{Berry--Keating--Connes} & \textbf{PT}\\
\midrule
Phase space & $(x, p) \in \Reals_+^{2}$ & $(u, p_u)$, $u = \log p$\\
Operator & $H = xp + 1/2$ & $\HPTBK = (u\,p_u + p_u\,u)/2$\\
Measure & $dx\,dp / 2\pi$ & $du\,dp_u / 2\pi$\\
Cut-off & $x \ge \ell_x$, $p \ge \ell_p$, $\ell_x \ell_p = 2\pi$ & $u_\star\,p_\star = 2\pi$ (PT Planck cell)\\
Regularisation & adelic (Connes) & atomic PT\\
Trace & $\Tr(K) = \sum_\rho h(\gamma_\rho)$ (Weil) & $\Tr_{\mathrm{reg}} = \sum_p \kappap [\text{Euler + coupling}]$\\
Holonomic coupling & none (zeta only) & $\kappap$ (bifurcation $\qplus/\qminus$)\\
$2\pi$ & not structurally interpreted & $[u, p_u] = i$, canonical Liouville\\
$1/8$ & corner Maslov phase (artefact) & $\sper^{2}/2 = c_1/4 = \sper/4$ (spinor cohomology)\\
Spin $\sper = 1/2$ & postulated & overdetermined by $\sper^{2}/2 = \sper/4$ modulo $\Tone$\\
\bottomrule
\end{tabular}
\caption{Berry--Keating--Connes vs PT comparison: the main
contribution of PT is intrinsic enrichment (no number is posited ad
hoc).}
\label{tab:comparaison_BKC}
\end{table}

\paragraph{Historical clarifications.}
The original operator of \citeauthor{BerryKeating1999} (1999) is
$H = xp$; the form $H = xp + 1/2$ appears in the later formalisation
of \citeauthor{Sierra2008} (2008) and Sierra--Townsend, which
regularises $xp$ via a Planck cell and exhibits the symmetric
ordering. The regularisation of \citeauthor{Connes1996}
(\citeyear{Connes1996}, \citeyear{Connes1999}) is of a different
nature: it operates on the adelic space of id\`ele classes, by a
non-commutative trace formula. The PT formulation is atomic on the
primes: neither adelic nor continuous on $\Reals_+$.

\subsection{Perspectives}
\label{sec:perspectives}

Three natural extensions.

\begin{itemize}[nosep,leftmargin=*]
\item \textbf{Other $L$-functions.} The Dirichlet prediction having
been falsified, the extension to $L$-functions must be rethought.
The cohomology of the spinor bundle under \emph{twist} by $\chi$
completely modifies $c_1$; a proper cohomological reformulation is
open.
\item \textbf{Maslov phase in continuous Berry phase.}
The hypothesis proposes that $1/(8q)$ exists at the level of the
continuous Berry phase along the spectral flow, but does not
manifest in the zero count. To be tested.
\item \textbf{Strict APS rigour.} Develop the formal index theorem
on the cornered manifold of the double barrier (strict K3). Not
urgent for the operatorial framework, important for structural
coherence.
\end{itemize}
